\documentclass[10pt,letterpaper,twocolumn]{article}

\usepackage[margin=0.7in]{geometry}
\usepackage[T1]{fontenc}
\usepackage{lmodern}
\usepackage{amsmath,amssymb}
\usepackage{graphicx}
\usepackage{float}
\usepackage{subcaption}
\usepackage{booktabs}
\usepackage{cite}
\usepackage[bookmarks=false]{hyperref}
\usepackage[font=small,labelfont=bf]{caption}

\setlength{\columnsep}{0.25in}
\setlength{\parskip}{0.25em}
\setlength{\parindent}{1em}
\renewcommand{\arraystretch}{1.05}
\setlength{\textfloatsep}{8pt plus 2pt minus 2pt}
\setlength{\floatsep}{6pt plus 2pt minus 2pt}
\setlength{\intextsep}{6pt plus 2pt minus 2pt}

\hypersetup{
    colorlinks=true,
    linkcolor=blue,
    citecolor=black,
    urlcolor=cyan
}

\title{Integrating DNA Sequence with Reference Epigenomic and Conservation Context for Variant Prioritization by Predicted Methylation Sensitivity}
\author{Samuel Zhang}
\date{\today}

\begin{document}
\maketitle

\begin{abstract}
Genetic variants can perturb regulatory sequence features associated
with local DNA methylation, but mutation-specific effects are rarely
available in matched experimental data. We developed a multimodal model
that predicts CpG methylation from a centered DNA sequence and a
reference-context vector containing seven epigenomic tracks and two
target-base conservation values. A DNABERT-2 sequence tower and an
independently pretrained context tower were combined by gated fusion
and optimized for M-value regression and binary methylation-state
classification.

On a held-out TCGA-BRCA probe-level split, the sequence--context model
achieved a beta-value MAE of 0.1014 and ROC-AUC of 0.9666, outperforming
sequence-only and context-only baselines. The trained model was then
applied to an ascertained cohort of 757 GDC-derived somatic synonymous
variant--CpG pairs using matched wild-type and mutant inputs. Under the
GRCh38 reference orientation, \textit{MSRA}, \textit{DDC}, and
\textit{CHD5} were the highest-ranked full-cohort case studies;
\textit{CHD5} was also the strongest candidate among 83 pairs associated
with held-out test probes. Each case study exceeded its finite matched
synonymous reference set, but none remained significant after
cohort-wide false-discovery-rate correction. Their predicted directions
were stable across 50 Monte Carlo dropout passes, although the broader
ranking was only moderately stable and failed a reverse-complement
robustness check. The framework is therefore presented as an
exploratory procedure for variant nomination rather than a causal or
statistically confirmed variant-effect estimator.
\end{abstract}

\section{Introduction}

Single-nucleotide variants can influence phenotype through both
protein-coding and regulatory mechanisms. Individual substitutions may
change transcription-factor motifs, chromatin accessibility, local DNA
organization, RNA processing, or translation without altering the
encoded amino acid. Prioritizing such variants remains difficult
because most observed substitutions lack matched functional
measurements.

DNA methylation is a widely measured regulatory state associated with
chromatin organization, transcription, and cellular identity
\cite{jones2012functions,schubeler2015function}. Sequence variation may
alter local methylation by disrupting regulatory motifs or changing the
sequence environment of a CpG. A model that learns sequence-associated
methylation patterns can therefore provide a useful ranking signal,
provided that its outputs are interpreted as predictions rather than
measured mutation effects.

Recent deep learning models show that DNA sequence contains substantial
information for predicting methylation and other epigenomic states
\cite{zhou2015predicting,kelley2016quential,avsec2021effective}.
Sequence alone, however, does not directly encode the chromatin context
in which a locus is interpreted. In this study, DNABERT-2 sequence
representations are combined with reference epigenomic and target-base
conservation measurements to predict locus-level methylation. The
trained predictor is then evaluated on matched wild-type and mutant
sequences while the reference-context vector is held fixed.

The resulting score,
$\Delta\hat{\beta}=\hat{\beta}_{\mathrm{mut}}-
\hat{\beta}_{\mathrm{wt}}$, measures the sensitivity of the model to a
specified substitution. Because training uses locus-level methylation
rather than paired mutation-effect labels, this score is not assumed to
equal the true biological effect. Somatic synonymous SNVs provide a
focused application class because they preserve amino-acid sequence
while retaining possible regulatory effects \cite{supek2014synonymous}.
The study therefore asks two questions: whether reference context
improves held-out methylation prediction, and whether the trained model
can produce a calibrated and reproducible ranking of candidate
variants.

\section{Related Work}

\subsection{Synonymous Variants as Regulatory Candidates}

Synonymous substitutions were historically interpreted mainly through
codon usage, translational selection, and the redundancy of the genetic
code \cite{ikemura1981codon}. Because they do not alter the encoded
amino acid, they were often treated as a comparatively neutral class of
coding variation. Subsequent work has shown that this interpretation is
incomplete. Synonymous substitutions can influence splice-site
recognition, RNA secondary structure, transcript stability, translation
rate, cotranslational folding, and other processes that depend on the
nucleotide sequence rather than only the protein sequence
\cite{sauna2011understanding}.

Coding exons can also carry regulatory information at the DNA level.
Transcription-factor occupancy and other sequence-dependent interactions
may overlap protein-coding regions, so a synonymous change can preserve
the amino-acid sequence while modifying a regulatory motif or local DNA
organization \cite{stergachis2013exonic}. Specific synonymous variants
have been associated with altered methylation, CpG organization, or
allele-specific regulatory behavior in loci including \textit{MLH1},
\textit{TP53}, and \textit{ADAMTS13}
\cite{alvarez2025,blackburn2016,hall2022}. These examples motivate the
use of synonymous SNVs as a focused application class: direct
amino-acid substitution is removed, but multiple sequence-level
mechanisms remain possible.

This focus does not imply that every synonymous variant is regulatory,
or that methylation is the only relevant outcome. Rather, synonymous
variants provide a useful setting for testing whether a sequence model
can distinguish substitutions that induce unusually large predicted
responses from a much larger background of weaker changes. The same
counterfactual framework could later be extended to missense,
noncoding, or other single-nucleotide variant classes, but synonymous
variants offer a comparatively interpretable starting point.

\subsection{Deep Learning for Methylation and Regulatory Prediction}

Deep learning has substantially expanded the range of regulatory
properties that can be predicted from genomic sequence. DeepCpG
combines sequence-derived representations with neighboring CpG
measurements to predict single-cell methylation states
\cite{angermueller2017deepcpg}. CpGenie uses convolutional sequence
features to predict CpG methylation from local genomic context
\cite{zeng2017predicting}. More broadly, sequence-to-function models
such as Enformer learn long-range regulatory relationships for gene
expression and functional genomic tracks
\cite{avsec2021effective}, while DNABERT-2 provides pretrained genomic
representations that can be adapted to downstream prediction tasks
\cite{zhou2023dnabert}.

These approaches differ in biological setting, input length, target
construction, and evaluation design, so their published metrics are not
directly interchangeable. The present study therefore does not treat an
external model's reported performance as a controlled baseline.
Instead, the quantitative comparison is internal: sequence-only,
reference-context-only, and fused models use the same CpG loci, targets,
probe split, and output objectives. This design isolates whether the
reference-context representation contributes predictive information
beyond the local sequence under the evaluation used here.

The multimodal setting is motivated by a basic limitation of sequence
alone. Nucleotide sequence constrains regulatory potential, but the same
motif may be interpreted differently across chromatin environments.
Conversely, epigenomic tracks can summarize accessibility or histone
state without fully representing the sequence that gives rise to local
binding and methylation patterns. Combining a pretrained sequence
representation with reference epigenomic and conservation features
therefore offers a practical way to test whether the two information
sources are complementary. In this study the context features are shared
reference measurements rather than patient-matched assays, so the model
is best understood as integrating local sequence with a fixed breast
 epithelial and evolutionary context.

\subsection{In Silico Mutagenesis and Variant Prioritization}

Sequence-to-function predictors are commonly interrogated through
\textit{in silico} mutagenesis: a wild-type sequence is evaluated,
one nucleotide is changed, and the difference in model output is used
as a variant score. This paired construction controls all other input
positions and makes the resulting score easy to rank. However, a model
can predict baseline molecular states accurately while remaining
sensitive to small perturbations for reasons that do not correspond to
true biological effects. Predictive performance and variant-response
reliability are therefore related but distinct questions.

The present analysis combines several forms of evidence around the
counterfactual score. Matched observed synonymous controls place each
response in the context of comparable substitutions from the same
candidate cohort. Monte Carlo dropout measures whether effect direction
and high-ranking status persist under stochastic inference. A
reverse-complement diagnostic tests dependence on sequence orientation,
and mutation-overlapping JASPAR scoring asks whether the leading
substitutions also modify recognizable motif compatibility. None of
these analyses replaces experimental perturbation, but together they
provide a more informative prioritization assessment than deterministic
rank alone.


\section{Methodology}
\label{sec:methodology}

\subsection{Methylation Targets and Data Partitioning}
\label{sec:data}

DNA methylation labels were generated by
\texttt{build\_training\_data.py} from the locally archived matrix
\texttt{TCGA-BRCA.methylation450.tsv.gz} and the hg38 manifest
\texttt{HM450.hg38.manifest.tsv.gz}. Solid-tissue-normal columns were
selected from the matrix header by the TCGA sample-type code
\texttt{-11}. This selected 97 normal-sample columns from 893 total
sample columns in the archived matrix. Manifest probes were restricted
to chromosomes 1--22, X, and Y, required nonmissing hg38 coordinates, a
valid 5000-bp reference window, and a centered \texttt{CG}
dinucleotide.
A probe was retained only when at least one selected normal-sample beta
value was observed, and the surviving methylation targets were
inner-joined to the valid sequence windows.

This procedure produced 418,486 CpG loci spanning 24 chromosomes. Of
these, 220,992 had median beta values above 0.5 and were assigned to the
methylated class, while 197,494 were assigned to the unmethylated class.
Each learning example represented one CpG locus with a cohort-level
target rather than one patient--locus pair.

For a CpG observed in the $S=97$ selected normal-sample columns, the
continuous target was the median methylation fraction
\[
\beta=\operatorname{median}(\beta_1,\ldots,\beta_S),
\]
with missing sample-level measurements omitted. To avoid undefined
values under the logit-like transformation, beta values were clipped to
$[10^{-4},1-10^{-4}]$ and converted to M-values,
\[
M=\log_2\left(\frac{\widetilde\beta}
{1-\widetilde\beta}\right).
\]
The binary target was $y=1$ when the unclipped median beta value exceeded
0.5 and $y=0$ otherwise. The regression and classification heads were
therefore trained on related but distinct representations of the same
locus-level methylation state.

CpG probes were divided into stratified training, validation, and test
sets in an 80\%/10\%/10\% ratio using random seed 42. Both sequence and
reference-context inputs associated with a probe remained in the same
partition. The split supports comparison across models on a common set
of held-out probes; it was not designed as a chromosome-level transfer
test.

\subsection{Sequence and Reference-Context Inputs}

For each target CpG, the sequence input was the central 1000 bp of a
stored 5000-bp hg38 reference window. The target CpG occupied the center
of the extracted sequence. Each window was tokenized using the
DNABERT-2 tokenizer and padded or truncated to a maximum of 1000 tokens.
This representation captures local regulatory sequence while remaining
tractable for fine-tuning and repeated counterfactual inference.

The reference-context input contained nine measurements. Seven were
MCF-10A breast epithelial BigWig signals: ATAC-seq accessibility and
H3K4me3, H3K27ac, H3K27me3, H3K9me3, H3K36me3, and H3K4me1. For each
track, the mean signal was extracted over a 100-bp interval centered on
the target CpG. The remaining two features were PhyloP-100way
conservation values at the target cytosine and guanine. The resulting
vector was
\[
\mathbf e\in\mathbb R^9.
\]
During data generation, a BigWig query returned \texttt{0.0} when
the chromosome was unavailable, the interval was out of range, the
summary statistic was missing, or an extraction error occurred.
Remaining NaN values were replaced by zero at model input, and a
nine-dimensional finite-value mask was concatenated with the feature
vector:
\[
\mathbf e'=[\mathbf e;\mathbf m_{\mathrm{ctx}}]
\in\mathbb R^{18}.
\]
The mask therefore distinguishes NaNs retained in the generated tables,
but it cannot distinguish a true biological zero from an upstream
fallback zero. The context features were shared reference measurements
rather than sample-matched measurements for each TCGA specimen. Their
role was to supply a consistent local chromatin and conservation
background that could complement the nucleotide sequence.

\subsection{Model Architecture}
\label{sec:model-architecture}

The sequence-only, reference-context-only, and fusion models used the
same two prediction tasks and output-head structure. The final fusion
model contained a DNABERT-2 sequence tower and an independently
pretrained reference-context tower, both producing 768-dimensional
representations.

\paragraph{Sequence tower.}
DNABERT-2 generated contextual token embeddings
\[
\mathbf H_{\mathrm{seq}}=T(\mathbf s)
\in\mathbb R^{N\times768},
\]
where $N$ is the tokenized sequence length. A one-dimensional
convolution with kernel size 3 and ReLU activation was applied across
tokens:
\[
\widetilde{\mathbf H}_{\mathrm{seq}}
=\operatorname{ReLU}
\left(\operatorname{Conv1D}(\mathbf H_{\mathrm{seq}})\right).
\]
A learned scalar attention score was assigned to each non-padding token,
normalized by softmax, and used to form a weighted sum,
\[
\mathbf x_{\mathrm{seq}}
=\sum_{i=1}^{N}\alpha_i\widetilde{\mathbf h}_i
\in\mathbb R^{768}.
\]
The convolution supplies a local token-level transformation, while the
attention-pooling layer produces one fixed-dimensional sequence
representation for the CpG-centered window.

\paragraph{Reference-context tower.}
The 18-dimensional context-plus-mask input was processed by a multilayer
perceptron with dimensions
\[
18\rightarrow128\rightarrow256\rightarrow256.
\]
LayerNorm, GELU activation, and dropout of 0.2 followed the first two
linear layers. A final linear projection mapped the 256-dimensional
output to
\[
\mathbf x_{\mathrm{ctx}}\in\mathbb R^{768},
\]
matching the sequence representation dimension.

\paragraph{Gated fusion and prediction heads.}
The sequence and context representations were normalized independently
and concatenated. A gating network with dimensions
$1536\rightarrow128\rightarrow2$ produced two sigmoid gates,
$g_{\mathrm{seq}}$ and $g_{\mathrm{ctx}}$. The gates were independent
and were not constrained to sum to one. The fused representation was
\[
\mathbf z=
 g_{\mathrm{seq}}\widetilde{\mathbf x}_{\mathrm{seq}}
 +g_{\mathrm{ctx}}\widetilde{\mathbf x}_{\mathrm{ctx}}.
\]
Separate classification and M-value regression heads were then applied
to $\mathbf z$. Each head used a $768\rightarrow256\rightarrow1$
architecture with GELU activation and dropout of 0.2 after the hidden
layer.

\subsection{Training Procedure and Prediction Objectives}
\label{sec:training}

Training proceeded in three stages. First, the sequence-only model was
trained end to end with AdamW using learning rate $2\times10^{-5}$,
weight decay 0.01, batch size 4, and eight gradient-accumulation steps.
Training ran for at most 10 epochs with mixed precision, gradient
clipping at 1.0, and a linear schedule with 10\% warmup. The
sequence-only implementation also included an attention-magnitude term
with coefficient 0.01 in addition to the two prediction losses.

Second, the reference-context-only model was trained independently from
scratch. It used AdamW with learning rate $10^{-3}$, weight decay
$10^{-4}$, batch size 128, and at most 15 epochs. A OneCycle schedule
used the first 10\% of steps for warmup. This stage learned the context
encoder, its 768-dimensional projection, and the two output heads.

Third, the sequence and context checkpoints were loaded into the fusion
model. The DNABERT-2 encoder, sequence convolution and pooling layers,
context multilayer perceptron, and context projection were frozen. Only
the modality-specific LayerNorm modules, gating network, and two output
heads remained trainable. These parameters were optimized with AdamW at
$5\times10^{-4}$, weight decay $10^{-4}$, batch size 4, eight
accumulation steps, and at most 10 epochs, again using mixed precision,
gradient clipping, and 10\% linear warmup. All stages used random seed
42.

The classification head produced a logit $c$ and probability
$\hat p=\sigma(c)$. The regression head produced an M-value prediction
$\hat M$. The joint training loss was
\[
L=L_{\mathrm{BCE}}+L_{\mathrm{Huber}},
\]
where binary cross-entropy was applied to the methylation-state target
and Huber loss with $\delta=1.345$ was applied to the M-value target.
Predicted M-values were converted back to beta space by
\[
\hat\beta=\frac{2^{\hat M}}{1+2^{\hat M}}.
\]
Models were compared using MAE and RMSE in both M-value and beta space,
and ROC-AUC for binary methylation state. Each configuration was trained
once. The implemented checkpoint rule selected the lowest validation
MAE between the classification probability and continuous beta target,
whereas final continuous metrics and all counterfactual scores were
computed from the inverse-transformed regression output.

\subsection{Counterfactual Sensitivity and Matched Calibration}
\label{sec:counterfactual-inference}

For each candidate substitution, the trained fusion model was evaluated
on matched wild-type and mutant sequences. The two inputs differed only
at the tested nucleotide, and the reference-context vector was held
fixed. The resulting model-sensitivity score was
\[
\Delta\hat\beta=
\hat\beta_{\mathrm{mut}}-\hat\beta_{\mathrm{wt}}.
\]
Positive values indicate a higher mutant beta prediction, while negative
values indicate a lower prediction. Candidate variants were ranked by
$|\Delta\hat\beta|$. Because the model was trained on locus-level
methylation rather than paired perturbation labels, this quantity was
used as a ranking score rather than a measured mutation effect.

To calibrate unusually large responses, every candidate was compared
with observed synonymous variants from the same frozen cohort. Eight
matching tiers were evaluated in order: (1) same gene, SBS96 context,
CpG-effect class, and absolute mutation--probe distance within 25 bp;
(2) SBS96, CpG-effect class, and distance within 50 bp; (3) SBS96,
CpG-effect class, and distance within 100 bp; (4) SBS6, CpG-effect
class, and distance within 100 bp; (5) SBS6, CpG-effect class, and
distance within 250 bp; (6) transition/transversion class, CpG-effect
class, and distance within 250 bp; (7) CpG-effect class and distance
within 250 bp; and (8) all observed synonymous variants. Here distance
matching used the absolute difference between the target and control
distances from their paired CpGs.

The narrowest tier containing at least 20 controls was selected. The
target itself and duplicate variant--probe identifiers were excluded.
Control sets larger than 1000 were subsampled without replacement using
a variant-specific NumPy random seed formed from the first eight bytes
of the SHA-256 digest of the variant identifier plus base seed 42; the
selected indices were then sorted. For target $i$ with control set
$C_i$, the empirical magnitude $p$-value was
\[
p_i=
\frac{1+\sum_{j\in C_i}
\mathbf1\left[|\Delta\hat\beta_j|\ge
|\Delta\hat\beta_i|\right]}
{|C_i|+1}.
\]
Benjamini--Hochberg correction was applied across all evaluated
candidates. This analysis asks whether a response is extreme relative
to comparable observed synonymous model responses.

\subsection{Ranking Stability and Sequence Annotation}

Ranking stability was evaluated with 50 Monte Carlo dropout passes over
the same ordered set of 757 variant--CpG pairs. Dropout layers were
activated while the remaining modules stayed in evaluation mode. Within
each pass, the wild-type and mutant member of a pair used the same random
seed so that their difference was computed under a matched dropout
realization. Stability was summarized by Spearman correlation with the
deterministic ranking, top-ten overlap, candidate rank distributions,
sign consistency, and top-ten and top-twenty frequencies.

A deterministic reverse-complement diagnostic separately transformed
the wild-type and mutant sequences and the two ordered target-base
conservation features. It was used to measure sensitivity to sequence
orientation, not to redefine the primary ranking. Finally, 41-bp
windows centered on the three leading full-cohort variants were scanned
against 841 JASPAR2022 vertebrate CORE matrices on both strands.
Wild-type and mutant scores were compared at identical
mutation-overlapping placements; placements were retained when at least
one allele had score greater than 5, and the largest absolute score
change was reported for each case study.

\section{Experiments and Results}
\label{sec:results}

\subsection{Held-Out Methylation Prediction}

The sequence--context model was evaluated against sequence-only and
reference-context-only models on the common held-out probe partition.
All three models used the same target definitions, train--validation--test
split, and evaluation procedure, allowing the comparison to focus on the
information contributed by each input configuration.

\begin{table}[H]
\centering
\caption{\textbf{Held-out predictive performance.} Lower RMSE and MAE
and higher ROC-AUC indicate better performance.}
\label{tab:performance}
\small
\resizebox{\columnwidth}{!}{
\begin{tabular}{lccc}
\hline
Metric & Context only & Seq. only & Seq. + context \\
\hline
RMSE (M-value) & 1.9934 & 1.7504 & \textbf{1.5353} \\
MAE (M-value) & 1.4830 & 1.2859 & \textbf{1.1569} \\
RMSE ($\beta$) & 0.2219 & 0.1901 & \textbf{0.1613} \\
MAE ($\beta$) & 0.1398 & 0.1156 & \textbf{0.1014} \\
ROC-AUC & 0.9230 & 0.9473 & \textbf{0.9666} \\
\hline
\end{tabular}}
\end{table}

The fusion model produced the best value for every reported metric.
Relative to sequence only, M-value MAE decreased from 1.2859 to 1.1569,
beta-value RMSE decreased from 0.1901 to 0.1613, and beta-value MAE
decreased from 0.1156 to 0.1014. The ROC-AUC increased by 0.0193, from
0.9473 to 0.9666. These changes were observed for both the continuous
regression output and the binary methylation-state output rather than
being confined to one objective.

The context-only model retained substantial predictive signal, with a
ROC-AUC of 0.9230, but it was weaker than the sequence-only model on all
five metrics. The improvement of the fusion model over both
single-modality models therefore supports a narrow conclusion: under the
shared probe-level split, the reference-context features contain
information that complements the sequence representation. The
comparison does not establish that the particular gating network is the
only or optimal fusion strategy, but it does show that the final result
cannot be explained by the context tower alone.

\subsection{Construction of the Application Cohort}

Candidate simple somatic mutations were retrieved from the GDC
TCGA-BRCA API using an any-transcript
\texttt{synonymous\_variant} filter. The accompanying
\texttt{Solid Tissue Normal} condition selected cases whose metadata
included a normal specimen; it did not indicate that the somatic
mutation was detected in that specimen. The frozen construction
pipeline mapped the stored gene symbol to cBioPortal and retained rows
with an available methylation-profile value. Those values were not used
as model inputs, counterfactual targets, or validation outcomes. The
resulting set is therefore treated as a finite, ascertained application
cohort rather than a representative catalogue of all TCGA-BRCA
synonymous mutations.

All coordinates and sequence operations used GRCh38/hg38. Each variant
was paired with the nearest HumanMethylation450 CpG within 2,499 bp so
that the substitution occurred within the stored 5000-bp CpG-centered
window. Counterfactual inference used the central 1000 bp. The initial
table contained 1,263 variant--CpG rows. Among these, 783 produced
exactly one valid ACGT substitution inside the model input, whereas 480
produced no sequence difference because the paired mutation was outside
the central 1000-bp window. Removing 17 rows that failed the
centered-CpG requirement and nine duplicate variant--probe identifiers
left 757 unique eligible pairs.

The deterministic ranking, matched-reference analysis, and Monte Carlo
dropout analysis all used this same ordered set of 757 identifiers.
Mapping each candidate probe back to the predictive split assigned 606
pairs to training, 68 to validation, and 83 to test. These represented
596, 65, and 79 unique CpG probes, respectively.

\begin{table}[H]
\centering
\caption{\textbf{Predictive-partition membership of the application
cohort.} Multiple variants can be paired with the same CpG probe.}
\label{tab:candidate-split-membership}
\small
\begin{tabular}{lrr}
\hline
Partition & Candidate pairs & Unique CpG probes \\
\hline
Training   & 606 & 596 \\
Validation & 68  & 65  \\
Test       & 83  & 79  \\
\hline
Total      & 757 & 740 \\
\hline
\end{tabular}
\end{table}

The complete 757-pair ranking is consequently an application of the
trained predictor across a mixed-partition cohort. The 83 test-probe
pairs provide a narrower analysis in which the associated CpG targets
were not used during model fitting or checkpoint selection.

\subsection{Full-Cohort and Held-Out Candidate Rankings}

Under the fixed GRCh38 reference orientation, the three largest
absolute deterministic responses in the full cohort were associated
with \textit{MSRA} (chr8:g.10428264C$>$T; cg14264678), \textit{DDC}
(chr7:g.50543912G$>$A; cg05346287), and \textit{CHD5}
(chr1:g.6128067C$>$T; cg12135344). Their predicted sensitivities were
$+0.2427$, $-0.1816$, and $-0.1606$, respectively. The signs indicate
that the mutant input received a higher prediction for \textit{MSRA}
and lower predictions for \textit{DDC} and \textit{CHD5}.

The three probes belonged to different predictive partitions:
\textit{MSRA} was associated with validation, \textit{DDC} with
training, and \textit{CHD5} with test. This distinction separates the
full-cohort application ranking from the held-out-probe subset. When the
ranking was restricted to the 83 test-probe pairs, \textit{CHD5}
ranked first, followed by \textit{AHDC1}
($\Delta\hat\beta=-0.1021$; full-cohort rank 7) and \textit{LHX6}
($\Delta\hat\beta=+0.0833$; full-cohort rank 9).

The deterministic test-subset ordering alone did not make the three
candidates equally strong. Their candidate-specific dropout summaries
show a clear separation between \textit{CHD5} and the next two held-out
ranks.

\begin{table}[H]
\centering
\caption{\textbf{Stability of the three highest deterministic
held-out-probe candidates.} MC intervals and ranks summarize 50 matched
dropout passes over the full cohort.}
\label{tab:heldout-top3-stability}
\small
\resizebox{\columnwidth}{!}{
\begin{tabular}{lrrrrr}
\hline
Gene & $\Delta\hat\beta$ & MC interval & Sign & Median rank & Top-10 \\
\hline
\textit{CHD5}  & $-0.1606$ & $[-0.2804,-0.0791]$ & 100\% & 2.0  & 88\% \\
\textit{AHDC1} & $-0.1021$ & $[-0.1517,-0.0239]$ & 100\% & 17.5 & 38\% \\
\textit{LHX6}  & $+0.0833$ & $[-0.0728,+0.1807]$ & 80\%  & 11.5 & 48\% \\
\hline
\end{tabular}}
\end{table}

\textit{CHD5} retained its negative direction in all 50 passes, its
5th--95th percentile interval excluded zero, and it appeared in the top
ten in 88\% of passes. \textit{AHDC1} also retained its direction and
had an interval excluding zero, but its median rank fell to 17.5 and it
appeared in the top ten in only 38\% of passes. \textit{LHX6} showed
weaker evidence: its interval crossed zero and sign consistency was
80\%. \textit{CHD5} is therefore the principal held-out-probe case
study, while \textit{AHDC1} and \textit{LHX6} are reported as the next
deterministic test-subset ranks rather than equivalent leads.

\subsection{Matched-Reference Calibration}

The three highest full-cohort responses were compared with finite sets
of matched observed synonymous controls selected by the prespecified
hierarchy. Each case study exceeded every member of its selected
control set: 30 controls for \textit{MSRA}, 27 for \textit{DDC}, and
191 for \textit{CHD5}. This yielded nominal empirical $p$-values of
0.0323, 0.0357, and 0.0052.

\begin{table}[H]
\centering
\caption{\textbf{Matched calibration of the three highest-ranked
full-cohort case studies.} None remained significant after cohort-wide
Benjamini--Hochberg correction.}
\label{tab:matched-null-top3}
\small
\resizebox{\columnwidth}{!}{
\begin{tabular}{lcrrrr}
\hline
Gene & Split & $\Delta\hat\beta$ & Controls & $p$ & $q$ \\
\hline
\textit{MSRA} & Val.  & $+0.2427$ & 30  & 0.0323 & 0.9614 \\
\textit{DDC}  & Train & $-0.1816$ & 27  & 0.0357 & 0.9614 \\
\textit{CHD5} & Test  & $-0.1606$ & 191 & 0.0052 & 0.9614 \\
\hline
\end{tabular}}
\end{table}

The finite control counts determine the resolution of the empirical
tests. When no control is as extreme as a target, the smallest
attainable value is $1/(|C_i|+1)$, which gives the reported values for
all three case studies. After Benjamini--Hochberg correction across the
candidate cohort, each had $q=0.9614$. The matched analysis therefore
shows that these deterministic responses were unusually large relative
to their available reference sets, but it does not support a
cohort-wide statistically significant discovery.

\begin{figure*}[t]
    \centering
    \begin{subfigure}{0.32\textwidth}
        \centering
        \includegraphics[width=\linewidth]{matched_null_top1_MSRA.png}
        \caption{\textit{MSRA}}
    \end{subfigure}
    \hfill
    \begin{subfigure}{0.32\textwidth}
        \centering
        \includegraphics[width=\linewidth]{matched_null_top2_DDC.png}
        \caption{\textit{DDC}}
    \end{subfigure}
    \hfill
    \begin{subfigure}{0.32\textwidth}
        \centering
        \includegraphics[width=\linewidth]{matched_null_top3_CHD5.png}
        \caption{\textit{CHD5}}
    \end{subfigure}
    \caption{\textbf{Matched observed-synonymous reference
    distributions.} Histograms show the matched control responses and
    vertical lines mark the three highest-ranked full-cohort case
    studies. Each exceeded its finite matched set, but none remained
    significant after cohort-wide correction.}
    \label{fig:matched-null-top3}
\end{figure*}

\subsection{Monte Carlo Dropout Stability}

Across the complete cohort, the stochastic absolute-effect rankings
showed moderate agreement with the deterministic ranking. The median
Spearman correlation was 0.6678, with values from 0.6151 to 0.6984
across the 50 passes. Median top-ten overlap was 45\%, and individual
passes retained between three and seven members of the deterministic
top ten. These results indicate that the broader upper tail should be
viewed as a region of high model sensitivity rather than a fixed total
ordering.

The three full-cohort case studies were more stable than the top-ten set
as a whole. All retained their predicted directions in every pass, and
none of their empirical 5th--95th percentile intervals crossed zero.
\textit{MSRA}, \textit{DDC}, and \textit{CHD5} appeared in the top ten
in 94\%, 82\%, and 88\% of passes, and in the top twenty in 100\%, 98\%,
and 96\%, respectively.

\begin{table*}[t]
\centering
\caption{\textbf{Monte Carlo dropout stability of the three
full-cohort case studies.} Intervals are empirical 5th--95th
percentiles across 50 passes.}
\label{tab:mc-dropout-stability}
\resizebox{\textwidth}{!}{
\begin{tabular}{lrrrrrrr}
\hline
Gene &
Deterministic $\Delta\hat\beta$ &
MC median $\Delta\hat\beta$ &
MC interval &
Sign consistency &
Median rank &
Top-10 frequency &
Top-20 frequency \\
\hline
\textit{MSRA} & $+0.2427$ & $+0.2004$ & $[+0.1164,+0.2560]$ &
100\% & 2 & 94\% & 100\% \\
\textit{DDC} & $-0.1816$ & $-0.1248$ & $[-0.2505,-0.0798]$ &
100\% & 5 & 82\% & 98\% \\
\textit{CHD5} & $-0.1606$ & $-0.1884$ & $[-0.2804,-0.0791]$ &
100\% & 2 & 88\% & 96\% \\
\hline
\end{tabular}}
\end{table*}

The stochastic analysis therefore supports the persistence of these
three large reference-orientation responses under inference variation,
particularly their effect directions. It provides weaker support for
the exact ordering of candidates outside this small leading group.

\subsection{Orientation Sensitivity and Motif Concordance}

An exploratory reverse-complement diagnostic preserved none of the
forward-orientation top ten after consistently transforming both the
sequence and the ordered target-base conservation features. The
ranking therefore depends strongly on the sequence orientation used by
the trained architecture. All primary candidate results in this paper
refer to the fixed GRCh38 reference orientation used during training
and inference.

For complementary sequence annotation, 41-bp wild-type and mutant
windows were scanned against 841 JASPAR2022 vertebrate CORE matrices on
both strands. The analysis compared each allele at the same
mutation-overlapping placement and reported the largest absolute score
change passing the prespecified threshold.

\begin{table}[H]
\centering
\caption{\textbf{Largest mutation-overlapping JASPAR score changes for
the three full-cohort case studies.}}
\label{tab:jaspar-top3}
\small
\resizebox{\columnwidth}{!}{
\begin{tabular}{llrrr}
\hline
Gene & Matrix & WT score & Mutant score & Change \\
\hline
\textit{DDC}  & HOXA2 MA0900.2  & $-11.2084$ & $5.5890$  & $+16.7974$ \\
\textit{CHD5} & NPAS2 MA0626.1  & $8.0016$   & $-2.9635$ & $-10.9651$ \\
\textit{MSRA} & ZNF423 MA0116.1 & $5.8926$   & $0.8482$  & $-5.0444$ \\
\hline
\end{tabular}}
\end{table}

The \textit{DDC} substitution crossed the threshold from a strongly
unfavorable to a positive HOXA2 score, consistent with a predicted motif
gain. The \textit{CHD5} and \textit{MSRA} substitutions reduced NPAS2
and ZNF423 scores below the threshold, consistent with predicted motif
losses. These motif changes provide sequence-level concordance with the
fact that the substitutions alter model-relevant local sequence, but
position-weight-matrix scores do not directly measure transcription-
factor occupancy or methylation change.

\subsection{Reproducibility}

The repository includes scripts for data construction, sequence and
context pretraining, fusion training, held-out evaluation,
counterfactual ranking, matched calibration, dropout stability, and
JASPAR scanning. In particular, \texttt{build\_training\_data.py}
defines methylation-column selection, probe filtering, sequence
validation, and the train--validation--test split, while
\texttt{matched\_synonymous\_null.py} defines the complete control
hierarchy and deterministic subsampling rule. The deterministic,
calibration, and stability outputs were audited to contain the same 757
identifiers.

The recorded environment used Python 3.10.20, PyTorch 2.6.0 with CUDA
12.4, Transformers 5.8.0, NumPy 2.2.6, pandas 2.3.3, scikit-learn
1.7.2, SciPy 1.15.3, Biopython 1.87, and \texttt{pyjaspar} 4.0.0. The
clean analysis snapshot corresponds to Git commit
\texttt{273378e6b8704d52b3711c3b3a4f4266a88e0212}.

The 50-pass dropout sweep used a fixed ordered cohort and matched random
seeds for each wild-type--mutant pair, requiring approximately 10
minutes 26 seconds in the recorded analysis environment. Exact rerunning
also requires the external \texttt{hg38.fa},
\texttt{HM450.hg38.manifest.tsv.gz},
\texttt{TCGA-BRCA.methylation450.tsv.gz}, reference BigWig tracks,
generated split and candidate tables, and trained checkpoints. The
archived methylation matrix contained 893 sample columns, of which 97
were selected by the \texttt{-11} normal-sample code, and had SHA-256
checksum
\texttt{71f7a02dd9ff849f43e05c6e54a9b8266349c9697601dc0450c9dc30f47679db}.
The selected sample identifiers were recorded in
\texttt{tcga\_methylation\_input\_audit.json}. A complete release
should retain this audit file and record the source URL or accession and
checksum for every other large external input.


\section{Discussion}

The predictive comparison supports a focused conclusion: reference
context adds information beyond local sequence under the present
probe-level split. The fusion model improved every reported metric, but
this result does not isolate the gating mechanism from other possible
fusion strategies. The counterfactual analysis then demonstrates how a
trained methylation predictor can nominate variants while keeping
$\Delta\hat\beta$ distinct from an experimentally observed mutation
effect.

\textit{CHD5} is the strongest held-out-probe result because it ranked
first among the 83 test-probe pairs and remained stable under dropout.
\textit{MSRA} and \textit{DDC} are larger and stable full-cohort model
responses, but their probes belong to the validation and training
partitions. None of the three survived cohort-wide correction, and the
reverse-complement failure prevents treating them as strand-invariant
biological priorities. Their appropriate role is therefore as
reference-orientation-specific case studies for subsequent validation.

The main limitations are the absence of direct perturbation validation,
the mismatch between TCGA bulk-tissue labels and MCF-10A reference
features, nearby overlapping windows under the random probe split, one
training seed, and the checkpoint-selection mismatch. The frozen
candidate cohort is also shaped by cBioPortal methylation availability,
even though those values were not used downstream. In addition,
any-transcript synonymous eligibility combined with storing the first
returned gene symbol may affect gene-based matching for some cohort
members. These issues constrain the strength and scope of the candidate
claims but do not change the study's role as an exploratory
prioritization analysis.

\subsection{Future Work}

The highest-value next analyses are forward--reverse-complement averaged
inference, fixed-release transcript reannotation with repeated control
matching, and a dedicated test-only candidate analysis. Genomic-block
or chromosome-level evaluation and repeated training seeds would
provide stronger evidence of generalization. Ultimately, targeted
introduction of prioritized variants followed by locus-specific
methylation measurement would be required for biological validation.

\section{Conclusion}

The sequence--context model improved held-out methylation prediction
relative to both single-modality baselines and was used to rank 757
somatic synonymous variant--CpG pairs. \textit{MSRA}, \textit{DDC},
and \textit{CHD5} were stable full-cohort case studies, while
\textit{CHD5} was the principal held-out-probe result. None was
significant after cohort-wide correction, and the rankings were
orientation-sensitive. The framework therefore provides a reproducible
procedure for hypothesis generation and candidate nomination rather
than causal variant-effect discovery.

\bibliographystyle{ieeetr}
\bibliography{references}
\end{document}

